% Ad Familiares 15.2 --- headnote
% ad-familiares-15-02, 22 September 51 BC, camp at Cybistra

\textit{Cicero's second official dispatch from Cilicia
to the consuls, praetors, tribunes of the plebs, and
Senate, written from the camp at Cybistra on the 22nd
of September 51 BC (Perseus dateline: \textit{Scr. in
castris ad Cybistra vel xi vel x K. Oct. a. 703 (51)}).
The opening greeting --- ``If you are in good health,
it is well; I myself am in good health'' --- is the
formula \textit{si vos valetis bene est, ego quidem
valeo}, preserved here in its full abbreviated form
\textit{S. v. v. b. e. e. q. v.} The letter is the
companion to 15.1, sent at roughly the same moment, and
should be read alongside it. 15.1 reports the Parthian
crossing of the Euphrates; this one reports what Cicero
did about the corresponding instability in Cappadocia.}

\textit{The substantive thread is the relief of King
Ariobarzanes III of Cappadocia (Eusebes and
Philorhomaeus, ``Pious'' and ``Lover-of-Rome'') from a
plot against his life. The Senate, under a decree
moved by Cato, had commended the young king's safety to
Cicero --- ``a decree which had never been passed by
our order in respect of any king'' --- and Cicero, in
the three days his army paused at Cybistra, did the
deferred diplomatic work. He summoned Ariobarzanes,
walked him through the Senate's commendation, advised
him to look to his own safety; the king at first denied
any danger, then on the following day came back to
camp in tears with disclosures from his brother
Ariarathes that a plot was in fact afoot. Cicero
declined to lend him cavalry --- the army was needed at
the Cilician frontier, and the danger to Syria was
nearer --- but advised the young king on how to use his
royal authority to deal with the conspirators. The
closing sentence (``it appears with reason that you
have shown so great care and diligence for his safety'')
is precisely the formula of a proconsul reporting that
the Senate's vote has been vindicated by events: a
prudent piece of political bookkeeping at the moment
Cicero needs the Senate to think well of him.}
